\begin{mistake}{2026-06-15}{04}

\begin{problemcontent}
函数 \( f(x) = \ln |(x-1)(x-2)\cdots(x-n)| \) 的驻点个数为 \underline{\hspace{1cm}}.
\end{problemcontent}

\begin{solutioncontent}
\textbf{解法一（多项式与罗尔定理）}：
设 \( P(x) = (x-1)(x-2)\cdots(x-n) \)，它是具 \( n \) 个互异实根的 \( n \) 次多项式，根为 \( 1, 2, \dots, n \)。
原函数定义域为 \( x \neq 1, 2, \dots, n \)，求导得：
\[ f'(x) = \frac{P'(x)}{P(x)} \]
令 \( f'(x) = 0 \)，即 \( P'(x) = 0 \) 且 \( P(x) \neq 0 \)。
由罗尔定理，\( P(x) \) 的相邻两根之间至少包含 \( P'(x) \) 的一个根。\( n \) 个根构成了 \( n-1 \) 个开区间 \( (k, k+1) \)，故 \( P'(x) \) 至少有 \( n-1 \) 个实根。
又因 \( P'(x) \) 是 \( n-1 \) 次多项式，至多有 \( n-1 \) 个实根。因此它恰好有 \( n-1 \) 个不为整数的实根。
驻点个数即为 \( n-1 \)。

\textbf{解法二（单调性与零点定理）}：
利用对数性质展开并求导（注意 \( (\ln|x|)' = 1/x \)）：
\[ f'(x) = \sum_{k=1}^n \frac{1}{x-k} \]
令 \( g(x) = f'(x) \)，求导得：
\[ g'(x) = \sum_{k=1}^n -\frac{1}{(x-k)^2} < 0 \]
\( g(x) \) 在每个连续区间 \( (k, k+1) \) 内单调递减。因区间两端极限分别为 \( +\infty \) 和 \( -\infty \)，由零点定理，每个区间有且仅有 \( 1 \) 个根。
另外区间 \( (-\infty, 1) \) 上 \( g(x) < 0 \)，区间 \( (n, +\infty) \) 上 \( g(x) > 0 \)，无根。
共 \( n-1 \) 个区间，故有 \( n-1 \) 个驻点。
\end{solutioncontent}

\mistakeknowledge{
\begin{itemize}
  \item \textbf{核心公式}：\( (\ln |x|)' = \frac{1}{x} \)，无论 \( x \) 取正负，求导后均无绝对值。
  \item \textbf{多项式与罗尔定理}：\( n \) 次多项式若有 \( n \) 个互异实根，则其导数必有且仅有 \( n-1 \) 个互异实根。
  \item \textbf{零点定理结合单调性}：单调连续函数若在区间两端趋近于异号无穷，则区间内必有唯一零点。
\end{itemize}}

\end{mistake}
